\subsection{Contexto Histórico}
Desenvolvido principalmente por cientistas e engenheiros em laboratórios, o software anterior a década de 1970 era considerado por muitos como um bem público  \citep{levy1984hackers}. Sem grandes restrições, o ato de compartilhar, modificar e construir programas colaborativamente fazia parte da cultura de pesquisa dos laboratórios \citep{von2003open}. Tal comportamento, tornou-se ainda mais evidente em 1969, quando cientistas e engenheiros passaram a compartilhar programas de computador através da primeira rede transcontinental de computadores, conhecida como \textit{ARPANET}. Desenvolver software naquela época envolvia, em essência, o compartilhamento de informações.

Em meados da década de 1970,  Bill Gates, principal fundador da Microsoft, ao enfrentar problemas com a pirataria do seu sistema Altar BASIC, escreve uma "Carta Aberta aos Hobbistas"~(do inglês, \textit{Open Letter to Hobbysts}), onde critica o frequente compartilhamento de código fonte proprietário \citep{KAVANAGH20041}. A declaração de Gates, um tanto quanto oposta a cultura de compartilhamento dos laboratórios, se unia a opinião de outros importantes nomes da indústria de software, incluindo, entre eles, Larry Ellison, co-fundador da Oracle.  O grupo de desenvolvedores formado por pessoas como Gates e Ellison, argumentava principalmente sobre a importância de uma indústria de software sólida, desenvolvida a partir de projetos com código fonte restrito. 

Em pouco tempo, os argumentos que alavancavam a ideia de software proprietário não só se espalharam entre outras companhias de software, mas também representaram uma diminuição na cultura de compartilhamento de código, fazendo com que a maioria dos então "\textit{hackers} de laboratório"~aderissem por variadas razões, a ideia de código fonte proprietário \citep{4686552}. O crescente domínio de empresas enriquecidas sob o desenvolvimento de software com código fonte restrito passava então a ser evidente.

Em 1984, vivenciando o enfraquecimento da cultura de compartilhamento de código, Richard Stallman, pesquisador do Laboratório de Inteligência Artificial do MIT, totalmente oposto a ideias tais como as de Gates e Ellison, publica o \textit{GNU Manifesto}, documento onde introduz a definição de software livre (do inglês, \textit{free software})~\citep{stallman1985gnu}. De acordo com Stallman, o software ideal deveria respeitar a liberdade do usuário, incluindo a total e irrestrita permissão de executar, estudar, distribuir e modificar o código fonte dos programas~\citep{stallman2002free}. Entre as críticas do manifesto, estava a ideia de que estas liberdades não eram garantidas por programas com código fonte fechado e que, por esta razão, os ideais de software deveriam ser novamente reconsiderados. Nascia assim, o conceito até hoje conhecido como software livre.

O que garante que um software irá se manter livre é a sua licença \citep{KAVANAGH2004297}. Diferente do desenvolvimento proprietário, onde a lei de propriedade intelectual protege os direitos de um autor contra a apropriação indébita por terceiros, as licenças de código aberto são projetadas para garantir os direitos de futuros usuários contra qualquer apropriação indevida do programa~\citep{von2003open}. A licença MIT\footnote{\url{https://opensource.org/licenses/mit}}, por exemplo, é uma licença permissiva que garante aos usuários os direitos de usar, copiar, modificar, mesclar, publicar, distribuir, licenciar e vender cópias do software desenvolvido. Rails\footnote{\url{https://rubyonrails.org/}}, um \textit{framework} web escrito em Ruby, é um exemplo de software livre sob a licença MIT.


\subsection{Visibilidade e Motivações}
Embora o termo software livre ainda seja desconhecido por muitos usuários, os projetos que seguem tal filosofia tem mudado substancialmente o cenário de desenvolvimento de software~\citep{BITZER20061}. O \textit{Apache HTTP Server}, por exemplo, é um projeto de software livre já há anos conhecido entre desenvolvedores \textit{web}. Mantido por voluntários, o \textit{Apache HTTP Server} hospedou somente em Julho de 2019\footnote{\url{https://w3techs.com/technologies/details/ws-apache/all/all}}, 44.3$\%$ dos \textit{websites} ativos na internet naquele mês. O Mozilla Firefox\footnote{\url{https://www.mozilla.org/firefox/}}, um dos maiores navegadores \textit{web} do mundo, é outro exemplo de projeto de software livre desenvolvido voluntariamente, que também representa grande adesão de usuários no mercado de software.

Não apenas usuários e voluntários se veem interessados em software livre, mas diversas companhias também já se encontram alinhadas a esta mesma filosofia~\citep{pinto2018challenges}. A linguagem de programação Swift, por exemplo, até então desenvolvida por funcionários da Apple, é atualmente um exemplo de projeto de software livre de sucesso. Até mesmo a Microsoft, reconhecida por sua popularidade com projetos de código proprietário, também lançou em Novembro de 2015 uma versão aberta do seu editor de texto, o Visual Studio Code. Estas companhias parecem não só estarem interessadas em tornar disponíveis publicamente o código fonte de seus projetos, mas também já existem evidências de funcionários sendo pagos para trabalharem nestes projetos de código aberto~\citep{dias2018drives}.

As motivações para manter e dar suporte a projetos de software livre são muitas: Para desenvolvedores de software, projetos de software livre podem ser uma ótima maneira de interagir com outros desenvolvedores, aprender novas tecnologias e poder se distrair programando algo diferente~\citep{raymond1999cathedral,torvalds2001just}. Para as companhias, os projetos de software livre podem se tornar uma grande fonte de conhecimento e trabalho voluntário, os quais indiretamente beneficiam o desenvolvimento de seus produtos.

 Entre as principais dificuldades enfrentadas neste tipo de ecossistema, encontra-se a necessidade de atrair novos desenvolvedores aos projetos desenvolvidos.
 
O fato de tais comunidades não atraírem novos contribuidores o suficiente, usualmente implica na falta de condições necessárias para se manter o desenvolvimento do programa, levando, por consequência, muitos projetos a inatividade. A inatividade de projetos de software livre afeta o ecossistema de desenvolvimento de software de maneira geral, já que projetos são, de maneira geral construídos sob o código de demais projetos.
